\documentclass[11pt,a4paper]{article}
\usepackage[margin=2.5cm]{geometry}
\usepackage{booktabs}
\usepackage{amsmath}
\usepackage{graphicx}
\usepackage{hyperref}

\title{GLCM Feature-Level Voxel-Spacing Weighting: Physical Rationale and Synthetic Validation}
\author{}
\date{}

\begin{document}
\maketitle

\section*{Purpose}
This note documents a refinement to the voxel-spacing-aware GLCM implementation in the corrected modified PyRadiomics prototype. The aim is to preserve the native discrete GLCM definition for each offset direction while applying physical voxel-spacing information only at the angular aggregation step. This is a software-behaviour validation experiment, not a clinical validation experiment.

\section*{Previous implementation}
The previous corrected implementation used relative physical offset weighting, but it still combined directional GLCM matrices before feature computation:
\[
  P^{*}(i,j) = \sum_{a \in \mathcal{A}} W_a P_a(i,j),
\]
followed by feature extraction from the combined matrix:
\[
  F_{\mathrm{old}} = F(P^{*}).
\]
Although this activates anisotropy-aware directional weighting, it is physically less conservative for non-linear Haralick descriptors because the distribution is altered before the descriptor formula is applied.

\section*{Refined implementation}
The refined implementation keeps each directional co-occurrence matrix separate:
\[
  P_a(i,j), \qquad a \in \mathcal{A}.
\]
For an offset vector \(a=(a_z,a_y,a_x)\) and voxel spacing \(s=(\Delta_z,\Delta_y,\Delta_x)\), it computes:
\[
  d_{\mathrm{index}}(a) = \sqrt{a_z^2+a_y^2+a_x^2},
\]
\[
  d_{\mathrm{phys}}(a) = \sqrt{(a_z\Delta_z)^2+(a_y\Delta_y)^2+(a_x\Delta_x)^2},
\]
\[
  r_a = \frac{d_{\mathrm{phys}}(a)}{d_{\mathrm{index}}(a)}, \qquad
  w_a = \frac{1}{r_a+\epsilon}, \qquad
  W_a = \frac{w_a}{\sum_b w_b}.
\]
The descriptor is then computed per direction and aggregated at the feature level:
\[
  F_{\mathrm{new}} = \sum_{a \in \mathcal{A}} W_a F(P_a).
\]
This is the preferred interpretation because it preserves the native discrete topology and avoids constructing a synthetic averaged GLCM prior to non-linear feature evaluation.

\section*{Isotropic equivalence}
When voxel spacing is isotropic, the implementation bypasses voxel-spacing weighting and uses standard PyRadiomics angular handling:
\[
  \Delta_z=\Delta_y=\Delta_x \Rightarrow F_{\mathrm{voxelSpacing}} = F_{\mathrm{default}}.
\]
This guard is essential. Without it, even inverse-distance weighting can change isotropic images by assigning different weights to axial and diagonal offsets. The corrected implementation therefore treats \texttt{weightingNorm = voxelSpacing} as an anisotropy-specific correction, not as a general replacement for standard GLCM aggregation.

\section*{Implemented code path}
The change was applied in:
\[
\texttt{<REPOSITORY_ROOT>/librerias/radiomics\_modificado\_corregido/glcm.py}
\]
The GLCM matrix calculation now stores normalized voxel-spacing angular weights without collapsing \(P_a\). A helper function performs the final NaN-aware angular average. Standard PyRadiomics behaviour is preserved when \texttt{weightingNorm} is absent, and existing PyRadiomics distance-weighting modes retain their original matrix-level aggregation behaviour.

\section*{Synthetic validation}
The experiment used a deterministic anisotropic 3D synthetic texture with spacing \((1.0,1.0,5.0)\) and an isotropic sanity image with spacing \((1.0,1.0,1.0)\). Extraction was run in isolated virtual environments for original and corrected modified PyRadiomics.

\begin{table}[h]
\centering
\caption{Backward compatibility and isotropic sanity checks.}
\begin{tabular}{lrrrr}
\toprule
Comparison & Features & Mean abs. diff. & Max abs. diff. & Changed \\
\midrule
Original anisotropic vs corrected default anisotropic & 75 & 0.000000 & 0.000000 & 0 \\
Corrected default isotropic vs corrected voxelSpacing isotropic & 75 & 0.000000 & 0.000000 & 0 \\
\bottomrule
\end{tabular}
\end{table}

\begin{table}[h]
\centering
\caption{Activation of voxel-spacing-aware extraction under anisotropic spacing.}
\begin{tabular}{lrrrr}
\toprule
Family & Features & Changed & Median abs. diff. & Median rel. diff. \\
\midrule
GLCM & 24 & 24 & 0.095800 & 0.101046 \\
GLDM & 14 & 10 & 0.246221 & 0.437710 \\
GLRLM & 16 & 16 & 0.094053 & 0.139830 \\
GLSZM & 16 & 0 & 0.000000 & 0.000000 \\
NGTDM & 5 & 5 & 0.022053 & 0.335726 \\
\bottomrule
\end{tabular}
\end{table}

\section*{Interpretation}
The revised GLCM path behaves coherently in the controlled validation. Default extraction remains exactly backward compatible with original PyRadiomics. The isotropic sanity check is now exact, indicating that \texttt{weightingNorm = voxelSpacing} does not introduce an artificial semantic change when all voxel dimensions are equal. Under anisotropic spacing, all 24 GLCM features changed, showing that the spacing-aware execution path is active.

Physically, this refinement is preferable to matrix-level GLCM averaging because it preserves each native directional co-occurrence distribution and weights only the final directional descriptor contributions according to relative physical offset length. This supports the central methodological claim: voxel-spacing-aware extraction separates geometric correction from interpolation-induced signal modification while preserving the native voxel array.

\section*{Generated files}
The synthetic validation outputs are stored in:
\[
\texttt{<REPOSITORY_ROOT>/experiments/synthetic\_spacing\_validation}
\]
Key outputs include \texttt{comparison\_backward\_compatibility.csv}, \texttt{comparison\_voxelSpacing\_effect.csv}, \texttt{comparison\_isotropic\_sanity.csv}, and \texttt{report\_synthetic\_spacing\_validation.tex}.

\end{document}
